\documentclass[11pt,a4paper]{article}

% =========================================================================
%  RAPPORT GRAFANA — LOKI TAAF Integration
%  Theme : Monitoring / observabilité — palette orange
% =========================================================================

\usepackage[utf8]{inputenc}
\usepackage[T1]{fontenc}
\usepackage[french]{babel}
\usepackage{lmodern}
\usepackage[a4paper,margin=2.2cm,top=2.6cm,bottom=2.4cm]{geometry}
\usepackage{xcolor}
\usepackage{fontawesome5}
\usepackage{tikz}
\usetikzlibrary{shapes.geometric,positioning,calc,arrows.meta}
\usepackage{titlesec}
\usepackage[most]{tcolorbox}
\usepackage{listings}
\usepackage{fancyhdr}
\usepackage{booktabs}
\usepackage{tabularx}
\usepackage{array}
\usepackage{colortbl}
\usepackage{enumitem}
\usepackage{microtype}
\usepackage[bookmarks=true,colorlinks=true,linkcolor=accent,urlcolor=accent,citecolor=accent]{hyperref}

% --- Palette Grafana -----------------------------------------------------
\definecolor{primary}{HTML}{161719}
\definecolor{accent}{HTML}{F46800}
\definecolor{accentdark}{HTML}{B84F00}
\definecolor{light}{HTML}{FFF7F0}
\definecolor{midgray}{HTML}{6B7280}
\definecolor{codebg}{HTML}{1E1F22}
\definecolor{codefg}{HTML}{E6E6E6}
\definecolor{codekw}{HTML}{FF9F40}
\definecolor{codestr}{HTML}{B5CEA8}
\definecolor{codecmt}{HTML}{6A9955}

\titleformat{\section}{\sffamily\Large\bfseries\color{primary}}{\thesection}{0.8em}{}[\vspace{-0.4ex}{\color{accent}\rule{\linewidth}{1.2pt}}]
\titleformat{\subsection}{\sffamily\large\bfseries\color{accent}}{\thesubsection}{0.6em}{}
\titlespacing*{\section}{0pt}{1.6em}{0.8em}

\pagestyle{fancy}
\fancyhf{}
\renewcommand{\headrulewidth}{0pt}
\fancyhead[L]{\footnotesize\sffamily\color{midgray}\faChartLine~Loki / Grafana — Monitoring TAAF}
\fancyhead[R]{\footnotesize\sffamily\color{midgray}Honorine Kylian}
\fancyfoot[C]{\footnotesize\sffamily\color{midgray}\thepage}

\lstdefinestyle{sqlstyle}{
  backgroundcolor=\color{codebg},
  basicstyle=\ttfamily\footnotesize\color{codefg},
  keywordstyle=\color{codekw}\bfseries,
  stringstyle=\color{codestr},
  commentstyle=\color{codecmt}\itshape,
  showstringspaces=false, breaklines=true,
  frame=none, framesep=8pt, framerule=0pt,
  xleftmargin=12pt, xrightmargin=12pt,
  columns=fullflexible, inputencoding=utf8, extendedchars=true,
  literate={é}{{\'e}}1 {è}{{\`e}}1 {ê}{{\^e}}1 {à}{{\`a}}1 {ô}{{\^o}}1 {î}{{\^\i}}1 {ç}{{\c{c}}}1 {ù}{{\`u}}1 {â}{{\^a}}1 {É}{{\'E}}1 {È}{{\`E}}1 {Ê}{{\^E}}1 {À}{{\`A}}1 {Ô}{{\^O}}1 {Ç}{{\c{C}}}1 {Ù}{{\`U}}1 {Â}{{\^A}}1
}
\lstset{style=sqlstyle,language=SQL}

\newtcblisting{sqlbox}[1][]{
  listing only,
  enhanced jigsaw, breakable,
  colback=codebg, colframe=accent,
  arc=2pt, boxrule=0pt, leftrule=3pt,
  left=8pt, right=8pt, top=4pt, bottom=4pt,
  listing options={style=sqlstyle,language=SQL}, #1
}

\newtcolorbox{infobox}[1][]{
  enhanced, colback=light, colframe=accent,
  arc=2pt, boxrule=0pt, leftrule=3pt,
  left=10pt, right=10pt, top=6pt, bottom=6pt,
  fontupper=\small, #1
}
\newtcolorbox{resultbox}[1][]{
  enhanced, colback=accent!10, colframe=accent,
  arc=2pt, boxrule=0pt, leftrule=3pt,
  left=10pt, right=10pt, top=6pt, bottom=6pt,
  fontupper=\small, #1
}

\setlist[itemize]{leftmargin=*,itemsep=2pt,topsep=4pt}
\setlist[enumerate]{leftmargin=*,itemsep=2pt,topsep=4pt}

\begin{document}
\pagenumbering{gobble}

% ---- PAGE DE TITRE ------------------------------------------------------
\begin{tikzpicture}[remember picture, overlay]
  \fill[primary] (current page.north west) rectangle ([yshift=-9cm]current page.north east);
  \fill[accent] ([yshift=-9cm]current page.north west) rectangle ([yshift=-9.25cm]current page.north east);
  % Logo : dashboard / gauges
  \node[anchor=center] at ([yshift=-4.6cm]current page.north) {%
    \begin{tikzpicture}[scale=1.0]
      \fill[white,rounded corners=4pt] (-2.4,-1.7) rectangle (2.4,1.7);
      \fill[primary,rounded corners=4pt] (-2.4,1.1) rectangle (2.4,1.7);
      \fill[accent] (-2.15,1.4) circle (0.13);
      % Barres de monitoring
      \fill[accent] (-1.9,-1) rectangle (-1.4,-0.2);
      \fill[accent] (-1.1,-1) rectangle (-0.6,0.2);
      \fill[accent] (-0.3,-1) rectangle (0.2,0.6);
      \fill[accent] (0.5,-1) rectangle (1.0,-0.4);
      \fill[accent] (1.3,-1) rectangle (1.8,0.8);
      % Ligne de tendance
      \draw[white,line width=1.5pt] (-1.65,-0.2) -- (-0.85,0.2) -- (-0.05,0.6) -- (0.75,-0.4) -- (1.55,0.8);
      % Points
      \foreach \x/\y in {-1.65/-0.2,-0.85/0.2,-0.05/0.6,0.75/-0.4,1.55/0.8} \fill[white] (\x,\y) circle (0.06);
      \node[font=\ttfamily\bfseries\small,white] at (-1.6,1.4) {dashboard};
    \end{tikzpicture}};
  \node[anchor=north, white, font=\sffamily\Huge\bfseries] at ([yshift=-6.7cm]current page.north) {Loki Integration};
  \node[anchor=north, accent, font=\sffamily\Large] at ([yshift=-7.65cm]current page.north) {Logs \& Métriques TAAF};
  \node[anchor=north, white!75!primary, font=\sffamily\small] at ([yshift=-8.4cm]current page.north) {Grafana \textbullet{} Prometheus \textbullet{} cAdvisor};

  \node[anchor=south west, primary, font=\sffamily\small] at ([xshift=2.2cm,yshift=3cm]current page.south west) {\textbf{\textcolor{accent}{Auteur}}};
  \node[anchor=south west, primary, font=\sffamily\small] at ([xshift=2.2cm,yshift=2.5cm]current page.south west) {HONORINE Kylian};
  \node[anchor=south east, primary, font=\sffamily\small] at ([xshift=-2.2cm,yshift=3cm]current page.south east) {\textbf{\textcolor{accent}{Module}}};
  \node[anchor=south east, primary, font=\sffamily\small] at ([xshift=-2.2cm,yshift=2.5cm]current page.south east) {BUT3 R\&T — 2025/2026};
  \fill[accent] ([yshift=1.3cm]current page.south west) rectangle ([yshift=1.4cm]current page.south east);
  \node[anchor=south, midgray, font=\sffamily\footnotesize] at ([yshift=0.6cm]current page.south) {IUT Réseaux \& Télécommunications — Université de La Réunion};
\end{tikzpicture}
\newpage

\pagenumbering{arabic}
\setcounter{page}{1}

{\sffamily
\hypersetup{linkcolor=primary}
\renewcommand{\contentsname}{\textcolor{accent}{\faList~} Sommaire}
\tableofcontents
}
\vspace{0.4cm}

\section{Contexte}

\begin{infobox}[title={\faSatelliteDish~Mission}]
Concevoir un système d'observabilité pour une base scientifique des \textbf{Terres Australes et Antarctiques Françaises (TAAF)} : surveillance du personnel isolé, des communications satellites et des équipements critiques via dashboards Grafana alimentés par requêtes SQL.
\end{infobox}

L'environnement polaire impose des contraintes uniques : isolement géographique, communications intermittentes via satellite, équipements soumis à des conditions extrêmes. Le tableau de bord doit donc fournir une vue opérationnelle immédiate aux superviseurs.

\medskip
\noindent\textbf{Stack technique :}
\begin{itemize}
  \item \textbf{Grafana} \textemdash{} visualisation et alerting
  \item \textbf{Prometheus} \textemdash{} stockage time-series
  \item \textbf{cAdvisor} \textemdash{} métriques conteneurs Docker
  \item \textbf{PostgreSQL} \textemdash{} données métier (personnel, communications, équipements)
\end{itemize}

% =========================================================================
\section{Requêtes de validation}

\subsection{Q1 — Volumétrie de la table de communications}

\begin{sqlbox}
SELECT COUNT(*) AS total_lignes
FROM communications_satellite;
\end{sqlbox}

\textbf{Pourquoi :} \texttt{COUNT(*)} comptabilise toutes les lignes (y compris celles avec valeurs nulles), donnant le nombre total d'événements de communication enregistrés.

\subsection{Q2 — Communications réussies sur 7 jours}

\begin{sqlbox}
SELECT COUNT(*) AS comms_reussies_7j
FROM communications_satellite
WHERE statut = 'Réussie'
  AND timestamp_comm >= NOW() - INTERVAL '7 days';
\end{sqlbox}

\textbf{Pourquoi :} double filtre sur le statut et la fenêtre temporelle. Cette requête sert d'indicateur de fiabilité court-terme dans le dashboard.

\subsection{Q3 — Personnel présent}

\begin{sqlbox}
SELECT COUNT(*) AS personnel_present
FROM personnel_base
WHERE statut = 'Présent';
\end{sqlbox}

\textbf{Pourquoi :} compte les effectifs réellement opérationnels sur site (exclut congés, missions extérieures, etc.).

\subsection{Q4 — Qualité de signal moyenne (Toulouse, 30 jours)}

\begin{sqlbox}
SELECT ROUND(AVG(qualite_signal), 2) AS qualite_moyenne
FROM communications_satellite
WHERE station_receptrice = 'Toulouse'
  AND timestamp_comm >= NOW() - INTERVAL '30 days';
\end{sqlbox}

\textbf{Affichage Grafana :} panel \texttt{Stat} avec seuils colorés \textemdash{} vert au-dessus de 70, orange entre 40 et 70, rouge en dessous de 40.

\subsection{Q5 — Top 3 des types de communication par volume}

\begin{sqlbox}
SELECT type_communication, SUM(volume_donnees_mb) AS volume_total
FROM communications_satellite
GROUP BY type_communication
ORDER BY volume_total DESC
LIMIT 3;
\end{sqlbox}

\textbf{Affichage Grafana :} \texttt{Bar chart} horizontal pour comparaison visuelle directe.

\subsection{Q6 — Heure de pic d'activité}

\begin{sqlbox}
SELECT EXTRACT(HOUR FROM timestamp_comm) AS heure,
       COUNT(*) AS nb_transmissions
FROM communications_satellite
GROUP BY heure
ORDER BY nb_transmissions DESC;
\end{sqlbox}

\begin{resultbox}[title={\faClock~Résultat}]
\textbf{Pic à 7h} avec 16 transmissions, cohérent avec le début des opérations journalières dans la base.
\end{resultbox}

\subsection{Q7 — Personnel isolé — détection de risque}

\textbf{Cible :} identifier une personne sans communication depuis plus de 7 jours \textbf{ET} en mission depuis plus de 4 mois.

\begin{sqlbox}
SELECT personnel_id, nom, prenom, fonction, statut,
       date_arrivee, derniere_communication,
       EXTRACT(DAY FROM (NOW() - date_arrivee))           AS duree_mission_jours,
       EXTRACT(DAY FROM (NOW() - derniere_communication)) AS inactivite_jours
FROM personnel_base
WHERE derniere_communication < NOW() - INTERVAL '7 days'
  AND date_arrivee <= NOW() - INTERVAL '4 months'
ORDER BY inactivite_jours DESC, duree_mission_jours DESC
LIMIT 1;
\end{sqlbox}

\begin{resultbox}[title={\faExclamationTriangle~Cas critique identifié}]
\textbf{Julie MOREAU} — Biologiste marine
\begin{itemize}
  \item En mission depuis 345 jours (\(\sim\)11 mois)
  \item Sans communication depuis 257 jours (\(\sim\)8 mois)
  \item Statut : En congé
\end{itemize}
\end{resultbox}

\textbf{Analyse des risques opérationnels :}
\begin{itemize}
  \item \textbf{Isolement critique} : impossibilité de confirmer l'état de santé ou la progression de la mission.
  \item \textbf{Durée anormalement longue} : risques de fatigue extrême, psychologiques et physiques.
  \item \textbf{Impact scientifique} : perte potentielle de données de recherche.
  \item \textbf{Risque organisationnel} : planification de relève compromise.
  \item \textbf{Contexte TAAF} : isolement géographique aggravé par les conditions climatiques.
\end{itemize}

\subsection{Q8 — Équipements critiques en panne (\textgreater{}48h)}

Le diagnostic croise pannes \textbf{équipement} et taux d'\textbf{échec de communication} par station :

\begin{sqlbox}
WITH pannes AS (
  SELECT equipement_id, nom_equipement, batiment AS station,
         derniere_verification
  FROM equipements_critiques
  WHERE statut = 'En panne'
    AND derniere_verification < NOW() - INTERVAL '48 hours'
)
SELECT p.nom_equipement, p.station,
       COUNT(*) FILTER (WHERE c.statut_transmission <> 'Réussie') AS nb_echecs,
       COUNT(*) AS total_comms,
       ROUND(100.0 * COUNT(*) FILTER (WHERE c.statut_transmission <> 'Réussie')
             / NULLIF(COUNT(*), 0), 1) AS pct_echec
FROM pannes p
LEFT JOIN communications_satellite c
       ON c.station_receptrice = p.station
      AND c.timestamp_comm >= p.derniere_verification
GROUP BY p.nom_equipement, p.station
ORDER BY pct_echec DESC;
\end{sqlbox}

\textbf{Affichage Grafana :} panel \texttt{Table} priorisé par \texttt{pct\_echec} décroissant pour focaliser l'intervention sur les équipements à fort impact.

\subsection{Q9 — Dashboard d'alertes unifié}

Trois flux d'alertes différents sont normalisés en une seule table pour affichage centralisé :

\begin{sqlbox}
WITH a_mission AS (
  SELECT 'Mission > 90j'::text AS type_alerte,
         personnel_id::text   AS identifiant,
         nom || ' ' || prenom AS libelle,
         date_arrivee         AS reference_time,
         EXTRACT(DAY FROM (NOW() - date_arrivee))::int AS severite_score
  FROM personnel_base
  WHERE date_arrivee <= NOW() - INTERVAL '90 days'
),
a_comms AS (
  SELECT 'Comms > 24h'::text,
         personnel_id::text,
         nom || ' ' || prenom,
         derniere_communication,
         EXTRACT(HOUR FROM (NOW() - derniere_communication))::int
  FROM personnel_base
  WHERE derniere_communication < NOW() - INTERVAL '24 hours'
),
a_eq AS (
  SELECT 'Équip. en panne'::text,
         equipement_id::text,
         COALESCE(nom_equipement, station),
         debut_panne,
         EXTRACT(HOUR FROM (NOW() - debut_panne))::int
  FROM equipements_critiques
  WHERE etat = 'en panne'
)
SELECT * FROM a_mission
UNION ALL SELECT * FROM a_comms
UNION ALL SELECT * FROM a_eq
ORDER BY type_alerte, severite_score DESC, reference_time ASC;
\end{sqlbox}

\begin{infobox}[title={\faIcon{tachometer-alt}~Panel \og Alertes critiques \fg{}}]
Trois familles d'alertes regroupées sous une vue unique \textemdash{} le superviseur dispose en un coup d'\oe{}il du panorama opérationnel : personnel en mission longue, silences radio, équipements en panne.
\end{infobox}

\section{Synthèse technique}

\renewcommand{\arraystretch}{1.3}
\begin{tabularx}{\linewidth}{@{}>{\bfseries}p{1.2cm} l X@{}}
\toprule
\rowcolor{light} Niveau & Type panel & Cas d'usage \\ \midrule
Q1--Q3 & Stat       & Indicateurs simples comptés \\
Q4     & Stat seuil & Qualité signal avec code couleur \\
Q5     & Bar chart  & Comparaison volumes par catégorie \\
Q6     & Time series& Distribution temporelle d'activité \\
Q7     & Table      & Détection de cas critiques \\
Q8     & Table tri  & Priorisation interventions \\
Q9     & Alert list & Dashboard centralisé \\
\bottomrule
\end{tabularx}

\section{Conclusion}

Ce travail démontre l'usage de Grafana au-delà du simple monitoring système : croiser plusieurs sources métier pour produire de l'information actionnable. Les CTE et fonctions d'agrégation PostgreSQL permettent d'encoder dans une seule requête une logique d'alerte complexe \textemdash{} chaque panel Grafana devient alors un capteur de risque opérationnel.

\vspace{0.5cm}
\begin{center}
{\sffamily\itshape\color{midgray}\large
\og En zone polaire, une donnée silencieuse est une alerte. \fg{}}
\end{center}

\end{document}
